%% Final JRSI manuscript, rhetorically tightened and structurally reshaped
\documentclass[openacc]{rsproca_new}

\usepackage[T1]{fontenc}
\usepackage[utf8]{inputenc}
\usepackage{lmodern}
\usepackage{microtype}
\usepackage{graphicx}
\usepackage{booktabs}
\usepackage{longtable}
\usepackage{array}
\usepackage{amsmath,amssymb}
\usepackage{threeparttable}
\usepackage{multicol}
\usepackage{url}
\usepackage{csquotes}

\addbibresource{refs.bib}

\jname{rsif}
\Journal{J R Soc Interface\ }

\begin{document}

\title{Dynamic compliance in the human pelvic ring: sex-specific organisation of childbirth-relevant deformation pathways}

\author{Petr Heny\v{s}$^{1}$, Michal Kucha\v{r}$^{1,2}$, Alexander Mor\'{a}vek$^{2}$, Niels Hammer$^{3,4}$}

\address{$^{1}$Institute of New Technologies and Applied Informatics, Faculty of Mechatronics, Informatics and Interdisciplinary Studies, Technical University of Liberec, Liberec, Czech Republic\\
$^{2}$Department of Anatomy, Faculty of Medicine in Hradec Kr\'{a}lov\'{e}, Charles University, Hradec Kr\'{a}lov\'{e}, Czech Republic\\
$^{3}$Division of Macroscopic and Clinical Anatomy, Gottfried Schatz Research Center, Medical University of Graz, Graz, Austria\\
$^{4}$Division of Macroscopic and Clinical Anatomy, Gottfried Schatz Research Center, Medical University of Graz, Graz, Austria; Department of Orthopedic and Trauma Surgery, University of Leipzig, Leipzig, Germany; Division of Biomechatronics, Fraunhofer IWU, Dresden, Germany}

\subject{biomechanics, evolutionary anthropology, obstetrics}

\keywords{biomechanics, human pelvis, modal analysis, finite-element modelling, childbirth, allometry, sacroiliac joint}

\corres{Michal Kucha\v{r}\\
\email{kucharm@lfhk.cuni.cz}}

\begin{abstract}
The human pelvis must simultaneously support habitual bipedal loading and accommodate the rare but extreme mechanical demands of childbirth. This dual role is usually discussed in terms of sexual dimorphism and static canal dimensions, yet pelvic function depends on how a compliant ring system gains access to, and organises, specific deformation pathways. We analysed 278 anatomically standardised lumbopelvic finite-element models derived from clinical CT using correspondence-preserving modal decomposition, allometric normalisation, standardised locomotor and childbirth-motivated load cases, and sensitivity analyses of ligamentous and cartilaginous constraints. Pelvic size explained a substantial share of global stiffness variation, indicating that size is a major determinant of pathway accessibility. After adjustment for scale and age, however, sex differences persisted in selected modes, most prominently modes~2 and~9, showing that sex reorganises ring mechanics rather than merely scaling them. Bilateral and unilateral stance recruited a shared low-order support backbone in both sexes. Under simulated inlet opening, mode~9 emerged as the principal inlet-opening pathway, but men concentrated a larger fraction of the response in this single pathway whereas women distributed deformation more broadly across coordinated low-order modes. Sensitivity analyses localised regulation of mode~9 to the pubic symphysis, anterior sacroiliac ligaments and sacroiliac cartilage interfaces. We therefore interpret the female pelvis not as generally lax or evolutionarily compromised, but as a highly specified balanced compliant ring that preserves habitual support while retaining selective access to childbirth-relevant deformation during parturition.
\end{abstract}

\maketitle

\begin{multicols}{2}

\section{Introduction}

The human pelvis must satisfy two demanding functions at once. It transmits load between the trunk and lower limbs during habitual stance and locomotion, yet it also defines the bony boundary of the birth canal. This dual role has long encouraged a geometric interpretation of pelvic biology: selection for efficient bipedalism is taken to constrain pelvic form, whereas selection for parturition favours enlargement of the canal \cite{RosenbergTrevathan2002,Trevathan2015,Grunstra2023}. In that framing, pelvic dimensions become the main currency of function.

Geometry is essential, but it is not the whole mechanical problem. The pelvis is a closed load-bearing ring whose behaviour depends not only on bony dimensions but also on material distribution, sacroiliac and symphyseal interfaces, and ligamentous constraints. Recent work has accordingly moved beyond a one-dimensional reading of the obstetrical dilemma. Wider female pelves do not necessarily incur a locomotor penalty \cite{Warrener2015}; childbirth must be understood in relation to load transfer, pelvic-floor support and developmental robustness as well as canal size \cite{Pavlicev2020,Haeusler2021,Grunstra2019,Grunstra2023}; and human pelvic form varies substantially within and among populations \cite{BettiManica2018,Auerbach2018}, with important developmental and ontogenetic influences on adult morphology \cite{VerbruggenNowlan2017,Novak2020,KurkiDecrausaz2015}. These lines of evidence suggest that static pelvimetry alone cannot capture pelvic function.

What remains limited is a population-scale framework that describes pelvic mechanics directly. Subject-specific finite-element studies have shown that sacroiliac constraints, pubic symphysis stiffness and load direction materially affect pelvic behaviour \cite{Eichenseer2011,Hammer2013}, consistent with broader evidence that the sacroiliac joint is a specialised load-transfer articulation with small but functionally relevant motion \cite{Vleeming2012,Kiapour2019,ChoKwak2020}. Yet case-by-case simulations do not easily reveal which features of mechanical behaviour are shared across individuals, which vary with scale, and which are organised by sex.

Here we address that problem using correspondence-preserving modal analysis of 278 anatomically standardised human pelves. Modal decomposition treats the pelvis not as a rigid set of diameters but as a repertoire of deformation pathways with different stiffness costs. This shift allows two questions to be separated. First, which pathways are mechanically accessible? Second, once a pathway is available, how is response organised under a given load? That distinction is central to childbirth-related mechanics.

Our central proposition is that the female pelvis is best understood not as an evolutionarily compromised or generally loosened structure, but as a highly specified compliant ring. In this framework, pelvic size should primarily modulate the accessibility of deformation pathways, whereas sex should primarily modulate how the ring organises response once those pathways are engaged. We therefore test whether size explains a substantial share of stiffness variation, whether sex differences persist after adjustment for scale and age, whether locomotor loads recruit a shared low-order support backbone, and whether the principal inlet-opening pathway is governed chiefly by the pubic symphysis and sacroiliac constraints rather than by static canal dimensions alone. If so, pelvic function is better described as an architecture of dynamic compliance than as a simple geometric compromise.

\section{Material and methods}

\subsection{Cohort and correspondence-preserving framework}

The analysis used the BoneDat resource, a standardised database of 278 clinical lumbopelvic CT scans from a Central European cohort spanning adolescence to old age \cite{HenysKuchar2025}. In the present derived workbooks, the cohort comprised 150 female and 128 male individuals, aged 16--91 years. BoneDat provides segmented anatomy, a shape-normalisation pipeline, a common reference surface and volumetric meshes, allowing each subject to be represented on a shared anatomical domain while preserving subject-specific morphology through deformation fields. This correspondence-preserving design is essential here because it permits direct mode-wise and region-wise comparison across the full cohort.

The original CT data were acquired on a Siemens Definition AS+ scanner and passed the technical quality control reported for BoneDat. Because the scans were acquired without a calibration phantom, mapped material properties are interpreted here as cohort-wise relative descriptors rather than direct clinical densitometry; phantom-less internal calibration was used to convert the released HU values to hydroxyapatite-equivalent density before elastic property assignment.

\subsection{Allometry and static morphology}

Allometric descriptors were taken from \texttt{allometry.xlsx}. The principal size variable was defined as the cube root of total pelvic volume relative to a template reference volume, which provides an appropriate whole-pelvis scale measure for a three-dimensional structure. Additional descriptors included total surface area, total mass, effective thickness and effective density. Static morphological dimensions were extracted from \texttt{anatomy\_data.xlsx}; the main variables were the anteroposterior inlet diameter (AP), pelvic inlet transverse diameter (PIT), subpubic angle, biacetabular width, biischiadic width, bituberous width, sacral width and iliopectineal eminence width. These variables were retained as the conventional pelvimetric descriptors against which the mechanical phenotype could be compared.

\subsection{Finite-element modelling and modal decomposition}

Each released deformation field was used to warp a tetrahedral template mesh into subject space, after which mechanics were evaluated on the common reference domain. Bone, sacroiliac cartilage and the pubic symphysis were represented on this shared mesh. Bone material properties were spatially heterogeneous and derived from mapped hydroxyapatite-equivalent density; sacroiliac cartilage and the pubic symphysis were assigned region-wise material properties. Linear isotropic elasticity was used throughout.

For each subject, pelvic mechanics were summarised through a generalised elastic eigenproblem,
\begin{equation}
K\phi_{k} = \lambda_{k} M\phi_{k}, \qquad k=1,2,\ldots,
\end{equation}
where $K$ and $M$ are the subject-specific stiffness and mass matrices on the mapped domain. Lower eigenvalues were interpreted as mechanically more accessible deformation pathways. We focused on modes~1--10, which the supplied documentation identifies as the biologically relevant low-order spectrum.

We distinguish two features of pelvic behaviour. \emph{Pathway accessibility} refers to where a mode lies in the stiffness spectrum: low-order, low-eigenvalue modes are easier to recruit mechanically. \emph{Pathway dominance} refers to how strongly a given load case concentrates response energy into that mode once the pathway has been engaged. This distinction is central to the interpretation of mode~9.

\subsection{Load cases and modal participation}

Static responses were computed for five standardised loading scenarios. Two represented habitual support mechanics: bilateral stance (\texttt{SP2leg}) and unilateral stance (\texttt{SP1leg}). Three represented childbirth-motivated probes: inlet opening (\texttt{LAB\_phase1}), outlet opening (\texttt{LAB\_phase2}) and terminal descent (\texttt{LAB\_phase3}). These scenarios were not intended as literal reproductions of labour, but as controlled mechanical probes that distinguish support-related from childbirth-relevant deformation pathways.

Static displacement fields were projected onto the modal basis by solving
\begin{equation}
(\Phi^{T}M\Phi)q = \Phi^{T}Mu,
\end{equation}
where $u$ is the static response, $\Phi=[\phi_{1},\ldots,\phi_{m}]$ is the modal basis, and $q$ is the vector of modal participation factors. Modal strain energies were then computed as
\begin{equation}
U_{i}=\frac{1}{2}\lambda_{i}q_{i}^{2}\|\phi_{i}\|_{M}^{2}, \qquad
\eta_{i}=\frac{U_{i}}{\sum_{j}U_{j}}.
\end{equation}
For interpretation and table construction, fractions were normalised over modes~1--10 so that the load-case summaries remained focused on the mechanically meaningful low-order spectrum.

\subsection{Sensitivity analysis}

To identify which anatomical constraints regulate the most relevant deformation pathways, we analysed workbook-based sensitivities of eigenvalues to ligament stiffness and cartilage-related modulus changes. For a ligament bundle with axial stiffness $EA_{\ell}$, the first-order modal sensitivity was
\begin{equation}
\frac{\partial \lambda_{i}}{\partial EA_{\ell}}
=
\frac{\phi_{i}^{T}(\partial K/\partial EA_{\ell})\phi_{i}}
{\phi_{i}^{T}M\phi_{i}}.
\end{equation}
For cartilage and symphyseal regions, the code reported a normalised energy-based sensitivity proxy,
\begin{equation}
S_{i,R}=\frac{\mathcal{E}_{i}(R)}{E_{R}\,\phi_{i}^{T}M\phi_{i}},
\end{equation}
which is proportional under linearisation to $\partial \lambda_{i}/\partial E_{R}$. These metrics were used to rank the importance of ligament groups, sacroiliac cartilage interfaces and the pubic symphysis for individual modes. Our primary mechanistic focus was mode~9, with modes~2 and~5 retained as support-related and outlet-related reference modes, respectively.

\subsection{Statistical analysis}

Static sex differences were assessed with Welch two-sample $t$-tests. Mode-wise sex effects were estimated using linear models of the form
\begin{equation}
\log(\lambda) \sim \log(\mathrm{scale}) + \mathrm{sex} + \mathrm{age}.
\end{equation}
The corresponding sex coefficient was interpreted as the percentage change in modal stiffness for males relative to females after adjustment for scale and age. To evaluate the limits of static pelvimetry, we also fitted female-only univariate models relating size-corrected eigenvalues of modes~9 and~2 to individual pelvic dimensions.

\section{Results}

\subsection{Cohort structure and static dimorphism}

The cohort comprised 278 individuals, including 150 women and 128 men, with no meaningful age difference between the sexes (table~\ref{tab:cohort}). As expected, female pelves were smaller in overall scale than male pelves ($0.925 \pm 0.045$ versus $0.998 \pm 0.045$, $p=6.51 \times 10^{-32}$), yet they showed larger AP and PIT dimensions, a wider subpubic angle, and greater biischiadic and bituberous widths. Sacral width showed no significant sex difference. The dataset therefore reproduces classic static dimorphism, but in a regionally patterned rather than uniformly expanded form.

\subsection{Size is a major mechanical determinant, but it does not remove sex structure}

Allometric regressions showed near-isometric scaling of total volume with pelvic scale (slope $=3.000$, $R^2=1.000$), shallower scaling of total surface area (slope $=1.684$, $R^2=0.894$), and increasing total mass with scale (slope $=2.679$, $R^2=0.902$), while effective density decreased modestly (slope $=-0.321$, $R^2=0.117$; table~\ref{tab:allometry}). Size therefore explains an important part of baseline stiffness variation and, by implication, the accessibility of deformation pathways.

After adjustment for scale and age, however, several modes retained clear sex differences (table~\ref{tab:modes}). The largest effects were observed in mode~2 and mode~9, where male pelves were stiffer than female pelves by $20.72\%$ and $20.07\%$, respectively. Significant differences were also present in modes~3, 4, 7, 8 and~10, whereas mode~1 remained effectively neutral. Thus, size correction reduces but does not eliminate sex-specific structure in the modal spectrum.

\subsection{Static pelvimetry captures only a limited share of mechanical variation}

In female-only models, PIT explained $13.5\%$ of the variation in size-scaled mode~9 eigenvalues and AP explained $9.2\%$; for mode~2, the strongest single predictor was PIT ($R^2=0.089$), with all other dimensions explaining less than $8\%$ of the variance (table~\ref{tab:staticassoc}). Static pelvimetry therefore retains some information about mechanical phenotype, but no single diameter captures more than a modest share of the relevant modal variation.

\subsection{Habitual support mechanics recruit a shared low-order backbone}

Under bilateral stance, both sexes showed an almost identical response dominated by mode~2 ($84.6\%$ in women and $84.5\%$ in men; table~\ref{tab:loadcases}). Unilateral stance broadened recruitment but preserved the same basic pattern: modes~1--3 accounted for most of the response in both sexes. Terminal descent similarly returned to a mode~2-dominated response ($77.2\%$ in women and $74.7\%$ in men), with mode~8 as the main secondary contributor. Across these support-weighted scenarios, sex differences in modal participation were small.

\subsection{Childbirth-related loading reveals sex-specific organisation of low-order pathways}

The inlet-opening scenario produced the clearest sex difference in response organisation. Mode~9 was the dominant response component in 265 of 278 subjects and therefore represents the principal inlet-opening pathway in the present cohort and modelling framework. Men placed a larger mean fraction of response energy in mode~9 than women ($52.5\%$ versus $46.4\%$), whereas women recruited substantially more mode~4 ($18.2\%$ versus $6.0\%$) and maintained a broader low-order distribution. Thus, mode~9 defines the main inlet-opening route, but women and men differ in how response is organised once that route becomes available.

Under outlet-opening loading, mode~5 became dominant in both sexes ($40.7\%$ in women and $39.2\%$ in men), consistent with an outlet-related pathway. Even in this scenario, however, women retained greater concurrent participation of mode~9 ($16.1\%$ versus $12.4\%$), whereas men showed greater involvement of mode~6. Childbirth-related loading therefore does not reveal a separate mechanical module so much as a reweighting of related low-order pathways.

\subsection{Mode~9 is governed by an anterior constraint system}

Sensitivity analysis identified the strongest mechanistic correlates of the principal inlet-opening pathway. At the interface level, mode~9 was most sensitive to the left and right sacroiliac cartilage interfaces and then to the pubic symphysis interface. At the ligament level, the largest effect was associated with the symphyseal ligament group, followed by the anterior sacroiliac ligaments and, at lower magnitude, the interosseous sacroiliac ligaments; posterior sacroiliac, sacrospinous and sacrotuberous ligaments were far less influential (table~\ref{tab:sensitivity}). Mode~9 is therefore regulated most strongly by an anteriorly weighted ring constraint system rather than by gross inlet width alone.

\section{Discussion}

Our central result is a separation between two mechanical layers of pelvic phenotype. Pelvic size explains a substantial share of overall stiffness and therefore sets much of the accessibility landscape for low-order deformation pathways. Sex, however, continues to structure selected modes after adjustment for scale and age, especially modes~2 and~9. The female pelvis is therefore not simply a smaller version of the male pelvis, nor does it appear globally loosened. Rather, it is mechanically specified in the way stiffness is distributed across the ring.

This distinction becomes most visible when static accessibility is separated from load-specific dominance. Mode~9 emerged as the principal inlet-opening pathway because it consistently occupied a childbirth-relevant position in the low-order spectrum and dominated the inlet-opening response in most subjects. Yet women did not show greater fractional concentration in mode~9 itself. Instead, they showed a broader low-order recruitment pattern once inlet-opening behaviour was engaged. Men, by contrast, concentrated more of the response into mode~9. This accessibility--dominance distinction is important because it avoids a misleading interpretation of the female pelvis as generally more lax. The difference lies in response organisation, not in diffuse softening of the whole ring.

The locomotor results reinforce that point. Bilateral stance, unilateral stance and terminal descent all relied primarily on a shared support-related modal backbone, with only modest sex differences in recruitment. Childbirth-related response is therefore layered onto a mechanical system that remains competent for habitual support. In that sense, the female pelvis can be read as a balanced compliant ring: it preserves everyday load transfer while retaining the ability to recruit childbirth-relevant pathways under the rare but biologically critical loads of parturition.

The regulation of mode~9 further clarifies why static pelvimetry is insufficient. Although PIT and AP carried some predictive signal, no single diameter explained more than a modest portion of the relevant modal variation. The principal inlet-opening pathway was governed instead by the pubic symphysis, anterior sacroiliac ligaments and sacroiliac cartilage interfaces, implicating the mechanics of the whole ring rather than canal size alone. This interpretation aligns with previous finite-element studies showing strong ligamentous effects on pelvic load transfer \cite{Eichenseer2011,Hammer2013} and with broader evidence that the sacroiliac joint allows small but functionally meaningful motion within a heavily constrained system \cite{Vleeming2012,Kiapour2019,ChoKwak2020}. It also refines how the pubic symphysis should be understood: not as a source of diffuse instability, but as part of a selectively regulated anterior mechanism relevant to childbirth \cite{Grunstra2019,Stolarczyk2021}.

These findings help reposition biomechanical interpretations of the human pelvis within the broader literature on the obstetrical dilemma. Recent syntheses have argued that childbirth cannot be reduced to a single trade-off between locomotion and canal size \cite{Pavlicev2020,Haeusler2021,Grunstra2023}, and experimental work shows that wider pelves need not impose a substantial locomotor penalty \cite{Warrener2015}. Our results are consistent with that view, but add a specific mechanical formulation: locomotor stability and childbirth-relevant compliance need not sit on one linear axis because they can be organised through partly distinct features of ring mechanics. From this perspective, the female pelvis is not well described as an evolutionarily compromised structure. It is better interpreted as a highly functional configuration that balances habitual support with conditional access to childbirth-relevant deformation.

The study remains limited to linear, small-deformation models and a Central European morphology cohort, and it does not incorporate pregnancy-specific soft-tissue changes, muscle forces, fetal interaction or obstetric outcomes. The omitted \texttt{pretension\_sensitivity.xlsx} workbook also means that our sensitivity interpretation is restricted to stiffness and interface effects. These limitations, however, define the next step rather than the present conclusion. A useful extension will be to test whether the same modal phenotype predicts behaviour in pregnancy-aware, nonlinear or outcome-linked models, and whether similar pathway structure is preserved across more diverse human populations.

\section{Conclusion}

In this population-scale analysis, pelvic size explained much of the global stiffness variation and thus much of the accessibility of deformation pathways, but it did not exhaust pelvic mechanics. After adjustment for scale and age, sex differences remained concentrated in selected modes, especially modes~2 and~9, while habitual support scenarios continued to rely on a shared low-order backbone. Mode~9 emerged as the principal inlet-opening pathway and was regulated chiefly by the pubic symphysis and anterior sacroiliac constraint system. Taken together, these findings support a positive interpretation of the female pelvis as a highly specified balanced compliant ring: not generally lax or evolutionarily compromised, but mechanically organised to preserve everyday support while retaining selective access to childbirth-relevant deformation.

\section*{Ethics}

This study used retrospectively acquired clinical CT data released through the BoneDat resource. According to the BoneDat dataset documentation, the study was approved by the Ethics Committee of University Hospital Hradec Kr\'{a}lov\'{e} (Ref: 202411I03P).

\section*{Data accessibility}

The source morphology dataset is described in BoneDat \cite{HenysKuchar2025}. Derived statistical tables, figure-generation scripts and supplementary processing code should be deposited in a public repository before submission, and the repository DOI inserted here.

\section*{Authors' contributions}

P.H. and M.K. designed the study. P.H. and M.K. performed the computational workflow. P.H., M.K. and A.M. analysed the data. P.H., M.K. and A.M. drafted the manuscript. N.H. contributed anatomical and biomechanical interpretation and critically revised the text. All authors approved the final version.

\section*{Competing interests}

The authors declare no competing interests.

\section*{Funding}

[Insert final funding statement and grant numbers.]

\section*{Acknowledgements}

We thank the BoneDat project for providing the standardised anatomical framework on which this analysis is based.

\end{multicols}
\begin{table}[p]
\centering
\caption{Cohort summary and static pelvic measures. Values are mean $\pm$ s.d. $p$-values are Welch two-sample $t$-tests (female versus male). AP, anteroposterior inlet diameter; PIT, pelvic inlet transverse diameter.}
\label{tab:cohort}
\begin{tabular}{lccc}
\toprule
Measure & Female $(n=150)$ & Male $(n=128)$ & $p$ \\
\midrule
Age (years) & $53.4 \pm 19.7$ & $54.5 \pm 18.9$ & 0.617 \\
Scale & $0.925 \pm 0.045$ & $0.998 \pm 0.045$ & $6.51 \times 10^{-32}$ \\
Effective density & $1.341 \pm 0.079$ & $1.321 \pm 0.074$ & 0.028 \\
AP (mm) & $121.0 \pm 10.0$ & $112.3 \pm 9.1$ & $4.63 \times 10^{-13}$ \\
PIT (mm) & $134.9 \pm 8.6$ & $127.7 \pm 7.6$ & $2.88 \times 10^{-12}$ \\
Subpubic angle (deg) & $85.8 \pm 6.1$ & $69.1 \pm 5.3$ & $7.30 \times 10^{-71}$ \\
Biacetabular width (mm) & $129.1 \pm 8.2$ & $122.0 \pm 7.2$ & $2.94 \times 10^{-13}$ \\
Biischiadic width (mm) & $109.9 \pm 7.1$ & $97.7 \pm 7.1$ & $9.40 \times 10^{-35}$ \\
Bituberous width (mm) & $136.5 \pm 8.3$ & $127.3 \pm 7.5$ & $3.03 \times 10^{-19}$ \\
Sacral width (mm) & $116.7 \pm 5.7$ & $117.5 \pm 6.5$ & 0.327 \\
Iliopectineal eminence width (mm) & $106.2 \pm 7.0$ & $97.9 \pm 6.0$ & $1.90 \times 10^{-22}$ \\
\bottomrule
\end{tabular}
\end{table}

\begin{table}[p]
\centering
\caption{Allometric scaling of geometric and material-derived quantities. Slopes are from log--log regressions against pelvic scale.}
\label{tab:allometry}
\begin{tabular}{lccc}
\toprule
Outcome versus scale & Slope & $p$ & $R^2$ \\
\midrule
Total volume & 3.000 & $<10^{-300}$ & 1.000 \\
Total surface & 1.684 & $2.53 \times 10^{-136}$ & 0.894 \\
Total mass & 2.679 & $2.48 \times 10^{-140}$ & 0.902 \\
Effective density & $-0.321$ & $5.07 \times 10^{-9}$ & 0.117 \\
\bottomrule
\end{tabular}
\end{table}

\begin{table}[p]
\centering
\caption{Mode-wise scale dependence and adjusted sex effects (modes~1--10). Correlations are Pearson's $r$ between $\log(\mathrm{scale})$ and $\log(\lambda)$ for unscaled and size-scaled eigenvalues. Adjusted sex effect is percent change (male relative to female) from $\log(\lambda) \sim \log(\mathrm{scale}) + \mathrm{sex} + \mathrm{age}$.}
\label{tab:modes}
\begin{tabular}{cccccc}
\toprule
Mode & $r$ (unscaled) & $p$ & $r$ (scaled) & $p$ & Sex effect \\
\midrule
1 & $-0.461$ & $4.96 \times 10^{-16}$ & 0.009 & 0.876 & $+0.85\%$ ($p=0.734$) \\
2 & $-0.331$ & $1.52 \times 10^{-8}$ & 0.374 & $1.20 \times 10^{-10}$ & $+20.72\%$ ($p=2.62 \times 10^{-18}$) \\
3 & $-0.467$ & $1.82 \times 10^{-16}$ & 0.233 & $8.84 \times 10^{-5}$ & $+11.23\%$ ($p=1.80 \times 10^{-7}$) \\
4 & $-0.509$ & $1.04 \times 10^{-19}$ & 0.273 & $3.70 \times 10^{-6}$ & $+15.90\%$ ($p=4.81 \times 10^{-9}$) \\
5 & $-0.607$ & $2.16 \times 10^{-29}$ & $-0.043$ & 0.480 & $-4.18\%$ ($p=0.106$) \\
6 & $-0.540$ & $5.15 \times 10^{-22}$ & 0.071 & 0.239 & $+3.39\%$ ($p=0.097$) \\
7 & $-0.778$ & $1.44 \times 10^{-57}$ & 0.179 & $2.76 \times 10^{-3}$ & $+5.13\%$ ($p=9.70 \times 10^{-4}$) \\
8 & $-0.398$ & $2.13 \times 10^{-11}$ & 0.304 & $2.37 \times 10^{-7}$ & $+14.45\%$ ($p=4.88 \times 10^{-11}$) \\
9 & $-0.369$ & $2.06 \times 10^{-10}$ & 0.388 & $1.97 \times 10^{-11}$ & $+20.07\%$ ($p=8.15 \times 10^{-20}$) \\
10 & $-0.630$ & $4.28 \times 10^{-32}$ & 0.218 & $2.56 \times 10^{-4}$ & $+8.66\%$ ($p=1.43 \times 10^{-5}$) \\
\bottomrule
\end{tabular}
\end{table}

\begin{table}[p]
\centering
\caption{Female-only associations between static pelvic dimensions and size-scaled eigenvalues. Values are $R^2$ and $p$ for univariate linear models of $\log(\lambda_{\mathrm{scaled}})$ versus a single dimension.}
\label{tab:staticassoc}
\begin{tabular}{lcccc}
\toprule
& \multicolumn{2}{c}{Mode 9} & \multicolumn{2}{c}{Mode 2} \\
\cmidrule(lr){2-3}\cmidrule(lr){4-5}
Dimension & $R^2$ & $p$ & $R^2$ & $p$ \\
\midrule
PIT & 0.135 & $3.63 \times 10^{-6}$ & 0.089 & $2.11 \times 10^{-4}$ \\
AP & 0.092 & $1.59 \times 10^{-4}$ & 0.001 & 0.64 \\
Iliopectineal eminence width & 0.086 & $2.77 \times 10^{-4}$ & 0.026 & 0.050 \\
Sacral width & 0.077 & $5.96 \times 10^{-4}$ & 0.003 & 0.49 \\
Biacetabular width & 0.018 & 0.11 & 0.053 & $4.58 \times 10^{-3}$ \\
Biischiadic width & 0.012 & 0.19 & 0.060 & $2.55 \times 10^{-3}$ \\
Bituberous width & 0.013 & 0.17 & 0.078 & $5.57 \times 10^{-4}$ \\
Subpubic angle & 0.071 & $9.84 \times 10^{-4}$ & 0.015 & 0.14 \\
\bottomrule
\end{tabular}
\end{table}

\begin{table}[p]
\centering
\caption{Mean modal participation fractions (\%) by loading scenario after normalisation over modes~1--10.}
\label{tab:loadcases}
\begin{tabular}{llrrrrrrrrrr}
\toprule
Scenario & Sex & M1 & M2 & M3 & M4 & M5 & M6 & M7 & M8 & M9 & M10 \\
\midrule
Bilateral stance & Female & 0.4 & 84.6 & 0.1 & 7.3 & 4.0 & 1.0 & 0.1 & 1.3 & 1.1 & 0.1 \\
Bilateral stance & Male   & 0.6 & 84.5 & 0.1 & 4.9 & 4.2 & 1.6 & 0.2 & 2.9 & 0.8 & 0.1 \\
Unilateral stance & Female & 39.6 & 37.8 & 15.3 & 2.6 & 1.5 & 1.2 & 0.2 & 0.5 & 0.7 & 0.6 \\
Unilateral stance & Male   & 40.4 & 37.2 & 14.1 & 1.9 & 1.4 & 2.3 & 0.2 & 1.0 & 0.7 & 0.9 \\
Inlet opening & Female & 0.6 & 4.4 & 4.5 & 18.2 & 3.8 & 1.6 & 1.0 & 15.3 & 46.4 & 4.3 \\
Inlet opening & Male   & 0.7 & 3.1 & 5.1 & 6.0 & 3.4 & 1.5 & 1.1 & 18.1 & 52.5 & 8.5 \\
Outlet opening & Female & 0.4 & 2.6 & 0.4 & 15.7 & 40.7 & 17.8 & 0.9 & 3.8 & 16.1 & 1.6 \\
Outlet opening & Male   & 0.4 & 2.4 & 1.1 & 16.4 & 39.2 & 22.1 & 1.2 & 2.8 & 12.4 & 2.1 \\
Terminal descent & Female & 0.3 & 77.2 & 0.2 & 8.8 & 1.6 & 0.7 & 0.5 & 10.6 & 0.1 & 0.1 \\
Terminal descent & Male   & 0.5 & 74.7 & 0.2 & 4.8 & 1.8 & 0.7 & 0.7 & 16.4 & 0.1 & 0.1 \\
\bottomrule
\end{tabular}
\end{table}

\begin{table}[p]
\centering
\caption{Sensitivity hierarchy for mode~9, with reference values for modes~5 and~2. Values are cohort means.}
\label{tab:sensitivity}
\begin{tabular}{lccc}
\toprule
Structure & Mode 9 sensitivity & Mode 5 sensitivity & Mode 2 sensitivity \\
\midrule
Left SIJ cartilage interface & $3.32 \times 10^{-3}$ & $6.63 \times 10^{-4}$ & $2.24 \times 10^{-4}$ \\
Right SIJ cartilage interface & $2.69 \times 10^{-3}$ & $3.74 \times 10^{-4}$ & $1.97 \times 10^{-4}$ \\
Pubic symphysis interface & $2.07 \times 10^{-3}$ & $6.93 \times 10^{-4}$ & $9.45 \times 10^{-6}$ \\
Symphyseal ligament group & $1.09 \times 10^{-6}$ & $7.86 \times 10^{-8}$ & $2.12 \times 10^{-9}$ \\
Anterior SIJ ligaments (bilateral mean) & $4.64 \times 10^{-7}$ & $9.77 \times 10^{-8}$ & $2.59 \times 10^{-8}$ \\
Interosseous SIJ ligaments (bilateral mean) & $2.40 \times 10^{-7}$ & $4.31 \times 10^{-9}$ & $4.68 \times 10^{-11}$ \\
Posterior SIJ ligaments (bilateral mean) & $1.27 \times 10^{-8}$ & $3.47 \times 10^{-8}$ & $2.87 \times 10^{-9}$ \\
Sacrospinous ligaments & $6.57 \times 10^{-9}$ & $3.07 \times 10^{-9}$ & $5.85 \times 10^{-9}$ \\
Sacrotuberous ligaments & $2.37 \times 10^{-8}$ & $1.24 \times 10^{-9}$ & $6.49 \times 10^{-10}$ \\
\bottomrule
\end{tabular}
\end{table}

\begin{figure}[p]
\centering
\fbox{\parbox[c][7cm][c]{0.9\linewidth}{\centering Placeholder for Figure~1\\
Cohort-to-template workflow and correspondence-preserving finite-element pipeline.}}
\caption{Study workflow. Clinical CT-derived pelvic anatomies were mapped onto a common template, converted into subject-specific finite-element models, decomposed into elastic eigenmodes, and analysed through load-case projection and sensitivity analysis.}
\label{fig:workflow}
\end{figure}

\begin{figure}[p]
\centering
\fbox{\parbox[c][7cm][c]{0.9\linewidth}{\centering Placeholder for Figure~2\\
Size-corrected modal spectrum and adjusted sex effects.}}
\caption{Size-corrected modal spectrum. Modes~2 and~9 show the largest adjusted sex effects after control for pelvic scale and age, whereas mode~1 acts as a near-neutral reference.}
\label{fig:eigs}
\end{figure}

\begin{figure}[p]
\centering
\fbox{\parbox[c][8cm][c]{0.9\linewidth}{\centering Placeholder for Figure~3\\
Representative deformation maps of modes~1, 2, 5 and~9.}}
\caption{Representative low-order deformation pathways. Mode~1 captures a conserved unilateral support component, mode~2 the baseline support-related ring mode, mode~5 an outlet-related pathway, and mode~9 the principal candidate inlet-opening pathway.}
\label{fig:modeshapes}
\end{figure}

\begin{figure}[p]
\centering
\fbox{\parbox[c][8cm][c]{0.9\linewidth}{\centering Placeholder for Figure~4\\
Stacked-bar modal participation fractions by loading scenario.}}
\caption{Load-case modal recruitment. Habitual support scenarios are built on a shared low-order modal backbone in both sexes. Under simulated inlet opening, men concentrate more of the response in mode~9, whereas women distribute the response more broadly across coordinated low-order pathways.}
\label{fig:recruitment}
\end{figure}

\begin{figure}[p]
\centering
\fbox{\parbox[c][7cm][c]{0.9\linewidth}{\centering Placeholder for Figure~5\\
Mode~9 versus static pelvic dimensions in women.}}
\caption{Static pelvimetry captures only a modest part of inlet-opening mechanics. PIT and AP are the strongest single predictors of size-corrected mode~9 stiffness in women, yet no single dimension explains more than a modest fraction of the variance.}
\label{fig:pelvimetry}
\end{figure}

\begin{figure}[p]
\centering
\fbox{\parbox[c][7cm][c]{0.9\linewidth}{\centering Placeholder for Figure~6\\
Sensitivity hierarchy of mode~9 across interfaces and ligament groups.}}
\caption{Sensitivity hierarchy for the principal inlet-opening pathway. Mode~9 is governed primarily by the pubic symphysis and anteriorly weighted sacroiliac constraints, whereas posterior ligament groups are much less influential.}
\label{fig:sensitivity}
\end{figure}

\clearpage
\printbibliography

\end{document}
